\documentclass[12pt]{article}
\usepackage{amsmath}
\usepackage{amssymb}
\usepackage{graphicx}
\usepackage{hyperref}
\usepackage{listings}
\usepackage{xcolor}
\usepackage{geometry}
\usepackage{booktabs}
\usepackage{tabularx}

\geometry{a4paper, margin=1in}

\definecolor{codegreen}{rgb}{0,0.6,0}
\definecolor{codegray}{rgb}{0.5,0.5,0.5}
\definecolor{codepurple}{rgb}{0.58,0,0.82}
\definecolor{backcolour}{rgb}{0.95,0.95,0.92}

\lstdefinestyle{mystyle}{
    backgroundcolor=\color{backcolour},   
    commentstyle=\color{codegreen},
    keywordstyle=\color{magenta},
    numberstyle=\tiny\color{codegray},
    stringstyle=\color{codepurple},
    basicstyle=\ttfamily\footnotesize,
    breakatwhitespace=false,         
    breaklines=true,                 
    captionpos=b,                    
    keepspaces=true,                 
    numbers=left,                    
    numbersep=5pt,                  
    showspaces=false,                
    showstringspaces=false,
    showtabs=false,                  
    tabsize=2
}
\lstset{style=mystyle}

\title{Laboratory: The Pure Polynomial Architecture for QR Codes}
\author{Advanced Abstract Algebra}
\date{}

\begin{document}
\maketitle

\section*{Introduction}
In this laboratory, we will explore the profound power of Python's **Duck Typing** to unify $GF(p)$ and $GF(p^n)$ arithmetic. We will build a single, universal \lstinline|Polynomial| class that performs algebra using standard operators (\lstinline|+|, \lstinline|-|, \lstinline|*|, \lstinline|/|). Because it doesn't care \textit{what} type of objects it is operating on, we can feed it integers modulo $p$ (for $GF(p)$ math) \textit{or} elements of our extension field (for $GF(p^n)$ math). 

Most importantly, we will completely discard the "integer representation" trick and the "lookup table" trick used in hardware optimization. Mathematically, elements of an extension field \textit{are} polynomials modulo a primitive polynomial! 

\section*{Problem 1: Universal Field Elements (GaloisFieldElement)}

\subsubsection*{Mathematical part}
The field $GF(p)$ is obtained by performing arithmetic modulo a prime number. A fundamental theorem states that $\mathbb Z/p\mathbb Z$ is a field if and only if $p$ is prime. This guarantees that every nonzero element has a unique inverse.

\subsubsection*{Implementation part}
\textbf{Implement:} the \lstinline|__add__|, \lstinline|__sub__|, \lstinline|__mul__|, \lstinline|__truediv__|, and \lstinline|__pow__| methods in the \lstinline|GaloisFieldElement| class. It encapsulates modulo arithmetic over an arbitrary field. Because it dynamically accesses \lstinline|self.field.modulus|, it will effortlessly work for both prime integer fields \textit{and} polynomial extension fields!

\section*{Problem 2: The Universal Polynomial}

\subsubsection*{Mathematical part}
Polynomial algebra operates identically regardless of the underlying field, as long as the coefficients support addition, subtraction, multiplication, and division.

\subsubsection*{Implementation part}
\textbf{Implement:} \lstinline|__add__|, \lstinline|__sub__|, \lstinline|__mul__|, \lstinline|__divmod__| (using explicit Euclidean long division), and \lstinline|__call__| (Horner's Method) in the \lstinline|Polynomial| class. Because it delegates all math to the coefficients' magic methods, it works for \textit{any} field! 

\section*{Problem 3: Primitive Polynomial Search}

\subsubsection*{Mathematical part}
Reed--Solomon codes require a field whose nonzero elements can be generated by repeated powers of a single element $\alpha^k$. A primitive polynomial is an irreducible polynomial whose root acts as such a generator for the multiplicative group of order $p^n - 1$.

\subsubsection*{Implementation part}
\textbf{Implement:} the \lstinline|is_primitive| function. Because you built the Universal Polynomial, this is a beautifully simple function. You will wrap the coefficients in \lstinline|GaloisFieldElement|, pass them to \lstinline|Polynomial|, and use \lstinline|Polynomial.__mul__| and \lstinline|Polynomial.__divmod__| to compute the exponentiations!

\section*{Problem 4: Extension Fields (ExtensionField)}

\subsubsection*{Mathematical part}
An extension field is created by adjoining a root of an irreducible polynomial. In hardware, this is optimized using Exponential and Logarithm lookup tables. **We will not do this.** 

Instead, an element of $GF(p^n)$ is \textit{literally} just a polynomial! To multiply two elements, you simply multiply their polynomials, and take the remainder modulo the primitive polynomial.

\subsubsection*{Implementation part}
\textbf{Implement:} Nothing! Because you correctly implemented \lstinline|GaloisFieldElement| to rely purely on \lstinline|field.modulus| in Problem 1, and the \lstinline|Polynomial| in Problem 2 supports the modulo operator, \lstinline|GaloisFieldElement| natively computes extension field arithmetic without any additional code. 

Marvel at how Python's duck typing just solved an entire laboratory level for you!

\section*{Problem 5: Reed-Solomon Encoding}

\subsubsection*{Mathematical part}
Reed--Solomon codes are evaluation codes. The generator polynomial $g(x)=\prod_{i=0}^{2t-1}(x-\alpha^i)$ forces every valid codeword to vanish at predetermined field elements. 

\subsubsection*{Implementation part}
\textbf{Implement:} \lstinline|get_generator_poly(num_ec_bytes, field)|. Notice that you just feed \lstinline|GaloisFieldElement| objects into the \textit{exact same} universal \lstinline|Polynomial| class!

\section*{Problem 6: Decoding I -- Syndromes}

\subsubsection*{Implementation part}
\textbf{Implement:} \lstinline|calculate_syndromes(message_poly, num_ec, field)|. Syndrome calculation uses \lstinline|Polynomial.__call__|!

\section*{Problem 7: Decoding II -- Linear System (PGZ)}

\subsubsection*{Implementation part}
\textbf{Implement:} \lstinline|pgz_error_locator(syndromes, field)| to construct the syndrome matrix and solve for the error locator polynomial. We have provided a fully generic \lstinline|solve_linear| utility that performs Gaussian Elimination. Because you implemented \lstinline|__bool__| on your elements, this utility doesn't even need to know what field it's operating on! It effortlessly solves the system over $GF(p^n)$ by purely relying on your object's mathematical overloads.

\section*{Problem 8: Decoding III -- Chien Search \& Linear Error Magnitudes}

\subsubsection*{Implementation part}
\textbf{Implement:} \lstinline|chien_search| to find the integer indices of the errors, and \lstinline|linear_error_magnitudes| to construct the Vandermonde matrix and calculate exact error magnitudes using \lstinline|solve_linear|.

\textbf{Test:} Run \lstinline|qr_code_demo.py| to test your unified OOP paradigm by generating a real QR code matrix!

\end{document}
